\documentclass[12pt]{article}
\usepackage[UTF8]{inputenc}
\usepackage{thaisupp}
\usepackage{enumitem}
\begin{document}
\title{ไฟฟ้ากระแสสลับ}
\maketitle
\section*{In-class Questions}
\begin{enumerate}[label=\arabic*.]
\item กระแสไฟฟ้าสลับ (AC) คือกระแสที่มีลักษณะอย่างไร \hfill \textbf{\small (Main)}
\begin{enumerate}[label=)]
\item ก) มีทิศทางเดียวตลอดเวลา
\item ข) เปลี่ยนทิศทางการไหลเป็นจังหวะสม่ำเสมอ
\item ค) มีค่าคงที่ตลอดเวลา
\item ง) ไหลในทิศทางเดียวแล้วกลับทิศทางช้าๆ
\end{enumerate}
\item ค่า Vrms หมายถึงอะไร \hfill \textbf{\small (Main)}
\begin{enumerate}[label=)]
\item ก) ค่าสูงสุดของแรงดันไฟฟ้า
\item ข) ค่าเฉลี่ยของแรงดันไฟฟ้า
\item ค) ค่ารากที่สองของค่าเฉลี่ยกำลังสองของแรงดัน
\item ง) ค่าแรงดันที่วัดได้จริงในวงจร DC
\end{enumerate}
\item แรงดันไฟฟ้าที่บ้านในประเทศไทยเป็น 220 V ค่า Vmax (ค่าสูงสุด) มีค่าเท่าใด \hfill \textbf{\small (Main)}
\begin{enumerate}[label=)]
\item ก) 220 V
\item ข) 311 V
\item ค) 440 V
\item ง) 110 V
\end{enumerate}
\item ในวงจร AC ที่มีตัวต้านทาน R ตัวเหนี่ยวนำ L และตัวเก็บประจุ C ต่ออนุกรมกัน อิมพีแดนซ์รวม Z มีค่าเท่าใด \hfill \textbf{\small (Main)}
\begin{enumerate}[label=)]
\item ก) Z = R + XL + XC
\item ข) Z = √[R² + (XL - XC)²]
\item ค) Z = R² + (XL - XC)²
\item ง) Z = √[(R + XL)² + XC²]
\end{enumerate}
\item กำลังไฟฟ้าเฉลี่ยในวงจร AC ขึ้นอยู่กับปัจจัยใด \hfill \textbf{\small (Main)}
\begin{enumerate}[label=)]
\item ก) เฉพาะ Vrms และ Irms
\item ข) Vrms, Irms และ cos φ
\item ค) เฉพาะความต้านทาน R
\item ง) เฉพาะความถี่ของกระแส
\end{enumerate}
\item ความถี่ของกระแสไฟฟ้าสลับในประเทศไทยเป็นเท่าไร \hfill \textbf{\small (Main)}
\begin{enumerate}[label=)]
\item ก) 60 Hz
\item ข) 50 Hz
\item ค) 100 Hz
\item ง) 40 Hz
\end{enumerate}
\item กระแสไฟฟ้าสลับมีค่าสูงสุด 10 A ค่า Irms มีค่าเท่าใด \hfill \textbf{\small (Main)}
\begin{enumerate}[label=)]
\item ก) 10 A
\item ข) 7.07 A
\item ค) 14.14 A
\item ง) 5 A
\end{enumerate}
\item วงจร AC มี R = 3 Ω, XL = 4 Ω และ XC = 0 Ω อิมพีแดนซ์ Z มีค่าเท่าใด \hfill \textbf{\small (Main)}
\begin{enumerate}[label=)]
\item ก) 3 Ω
\item ข) 4 Ω
\item ค) 5 Ω
\item ง) 7 Ω
\end{enumerate}
\item วงจร AC มี Vrms = 100 V, Irms = 5 A และ cos φ = 0.8 กำลังไฟฟ้าเฉลี่ยมีค่าเท่าใด \hfill \textbf{\small (Main)}
\begin{enumerate}[label=)]
\item ก) 500 W
\item ข) 400 W
\item ค) 250 W
\item ง) 625 W
\end{enumerate}
\item วงจร AC อนุกรมมี R = 8 Ω, XL = 15 Ω และ XC = 7 Ω อิมพีแดนซ์รวมและมุมเฟส φ มีค่าเท่าใด \hfill \textbf{\small (Main)}
\begin{enumerate}[label=)]
\item ก) Z = 10 Ω, tan φ = 1
\item ข) Z = 10 Ω, tan φ = 0.8
\item ค) Z = 17.2 Ω, tan φ = 1
\item ง) Z = 15.6 Ω, tan φ = 0.5
\end{enumerate}
\item ข้อใดไม่ใช่คุณสมบัติของกระแสไฟฟ้าสลับ (AC) \hfill \textbf{\small (Main)}
\begin{enumerate}[label=)]
\item ก) เปลี่ยนทิศทางการไหลเป็นจังหวะ
\item ข) มีความถี่ 50 Hz ในประเทศไทย
\item ค) สามารถแปลงระดับแรงดันได้ง่ายโดยใช้หม้อแปลง
\item ง) มีค่าคงที่ตลอดเวลาไม่เปลี่ยนแปลง
\end{enumerate}
\item เครื่องวัดไฟฟ้า AC ทั่วไปวัดค่าชนิดใด \hfill \textbf{\small (Main)}
\begin{enumerate}[label=)]
\item ก) ค่าสูงสุด (Peak value)
\item ข) ค่า RMS
\item ค) ค่าเฉลี่ย
\item ง) ค่าต่ำสุด
\end{enumerate}
\item วงจร AC มีความต้านทาน R = 6 Ω และมีการต่อตัวเหนี่ยวนำ L ในวงจร ถ้าอิมพีแดนซ์รวมเป็น 10 Ω ค่า XL มีค่าเท่าใด \hfill \textbf{\small (Main)}
\begin{enumerate}[label=)]
\item ก) 4 Ω
\item ข) 6 Ω
\item ค) 8 Ω
\item ง) 16 Ω
\end{enumerate}
\item ในวงจร AC ที่มีเฉพาะตัวเก็บประจุ กำลังไฟฟ้าเฉลี่ยเป็นเท่าใด \hfill \textbf{\small (Main)}
\begin{enumerate}[label=)]
\item ก) V × I
\item ข) V × I × cos φ
\item ค) 0
\item ง) V × I × sin φ
\end{enumerate}
\item สมการของกระแสไฟฟ้าสลับคือ I = I\_max sin(ωt) ข้อใดถูกต้องเกี่ยวกับ ω \hfill \textbf{\small (Main)}
\begin{enumerate}[label=)]
\item ก) ω คือความถี่เชิงมุม มีหน่วยเรเดียน/วินาที
\item ข) ω คือความถี่ มีหน่วยเฮิรตซ์
\item ค) ω คือมุมเฟสเริ่มต้น
\item ง) ω คือแรงดันไฟฟ้าสูงสุด
\end{enumerate}
\item วงจร AC อนุกรมมี R = 12 Ω, XL = 16 Ω และ XC = 10 Ω กระแสที่ไหลในวงจรเมื่อต่อกับแหล่งจ่าย 220 Vrms มีค่าเท่าใด \hfill \textbf{\small (Main)}
\begin{enumerate}[label=)]
\item ก) 22 A
\item ข) 11 A
\item ค) 10 A
\item ง) 5.5 A
\end{enumerate}
\item วงจร AC มี P = 1,000 W, Vrms = 200 V และ cos φ = 0.5 กระแส Irms มีค่าเท่าใด \hfill \textbf{\small (Main)}
\begin{enumerate}[label=)]
\item ก) 5 A
\item ข) 10 A
\item ค) 15 A
\item ง) 20 A
\end{enumerate}
\item ความสัมพันธ์ระหว่าง Vrms และ Vmax คือข้อใด \hfill \textbf{\small (Main)}
\begin{enumerate}[label=)]
\item ก) Vrms = Vmax × √2
\item ข) Vrms = Vmax / √2
\item ค) Vrms = Vmax × 2
\item ง) Vrms = Vmax / 2
\end{enumerate}
\item ในวงจร AC ที่มีเฉพาะตัวเหนี่ยวนำ L ความสัมพันธ์ระหว่างแรงดันและกระแสเป็นอย่างไร \hfill \textbf{\small (Main)}
\begin{enumerate}[label=)]
\item ก) กระแสนำหน้าแรงดัน 90°
\item ข) กระแสตามหลังแรงดัน 90°
\item ค) กระแสอยู่เฟสกับแรงดัน
\item ง) กระแสมีเฟสตรงข้ามกับแรงดัน
\end{enumerate}
\item ข้อใดไม่ใช่ข้อดีของไฟฟ้ากระแสสลับ (AC) เมื่อเทียบกับไฟฟ้ากระแสตรง (DC) \hfill \textbf{\small (Main)}
\begin{enumerate}[label=)]
\item ก) สามารถแปลงระดับแรงดันได้ง่ายโดยหม้อแปลง
\item ข) ส่งกำลังไฟฟ้าได้ระยะไกลด้วยแรงดันสูง
\item ค) ใช้งานง่ายในวงจรอิเล็กทรอนิกส์ดิจิทัล
\item ง) การสูญเสียพลังงานในสายส่งต่ำกว่า
\end{enumerate}
\item วงจร AC มีกำลังไฟฟ้าจริง (Real power) 500 W และกำลังไฟฟ้ารีแอกทีฟ (Reactive power) 500 VAR กำลังไฟฟ้าปรากฏ (Apparent power) มีค่าเท่าใด \hfill \textbf{\small (Main)}
\begin{enumerate}[label=)]
\item ก) 500 VA
\item ข) 707 VA
\item ค) 1,000 VA
\item ง) 250 VA
\end{enumerate}
\item ในวงจร AC ที่มีเฉพาะตัวเก็บประจุ C ค่าความต้านทานในการกันกระแส (Capacitive Reactance) XC มีค่าเท่าใด \hfill \textbf{\small (Main)}
\begin{enumerate}[label=)]
\item ก) XC = ωL
\item ข) XC = 1/ωC
\item ค) XC = ωC
\item ง) XC = R/ω
\end{enumerate}
\item ไฟฟ้ากระแสสลับในบ้านของประเทศไทยมีรูปคลื่นแบบใด \hfill \textbf{\small (Main)}
\begin{enumerate}[label=)]
\item ก) รูปคลื่นสี่เหลี่ยม
\item ข) รูปคลื่นฟันเลื่อย
\item ค) รูปคลื่นไซน์
\item ง) รูปคลื่นสามเหลี่ยม
\end{enumerate}
\item วงจร AC มี R = 40 Ω, XL = 30 Ω, XC = 10 Ω ต่อกับแหล่งจ่าย 220 Vrms กำลังไฟฟ้าที่ใช้ไปมีค่าเท่าใด \hfill \textbf{\small (Main)}
\begin{enumerate}[label=)]
\item ก) 1,210 W
\item ข) 2,420 W
\item ค) 605 W
\item ง) 968 W
\end{enumerate}
\item วงจรไฟฟ้าบ้าน 220 V มีแรงดันสูงสุดเท่าไร \hfill \textbf{\small (Main)}
\begin{enumerate}[label=)]
\item ก) 110 V
\item ข) 220 V
\item ค) 311 V
\item ง) 440 V
\end{enumerate}
\item คาบของกระแสไฟฟ้าสลับ 50 Hz มีค่าเท่าไร \hfill \textbf{\small (Extra)}
\begin{enumerate}[label=)]
\item ก) 0.02 วินาที
\item ข) 0.05 วินาที
\item ค) 0.01 วินาที
\item ง) 0.5 วินาที
\end{enumerate}
\item เครื่องวัดไฟฟ้าแสดงค่า 10 A สำหรับกระแส AC นี่คือค่าชนิดใด \hfill \textbf{\small (Extra)}
\begin{enumerate}[label=)]
\item ก) ค่าสูงสุด
\item ข) ค่า RMS
\item ค) ค่าเฉลี่ย
\item ง) ค่ากึ่งสูงสุด
\end{enumerate}
\item แรงดันไฟฟ้าสลับมีค่าสูงสุด 100 V ค่า Vrms มีค่าเท่าใด \hfill \textbf{\small (Extra)}
\begin{enumerate}[label=)]
\item ก) 50 V
\item ข) 70.7 V
\item ค) 141.4 V
\item ง) 100 V
\end{enumerate}
\item กระแสไฟฟ้าสลับมี Vrms = 150 V และ Irms = 4 A ค่ากำลังไฟฟ้าเฉลี่ยเป็นเท่าใด \hfill \textbf{\small (Extra)}
\begin{enumerate}[label=)]
\item ก) 150 W
\item ข) 600 W
\item ค) 424 W
\item ง) 300 W
\end{enumerate}
\item อิมพีแดนซ์ (Impedance) ในวงจร AC มีหน่วยเป็นอะไร \hfill \textbf{\small (Extra)}
\begin{enumerate}[label=)]
\item ก) โวลต์
\item ข) แอมป์
\item ค) โอห์ม
\item ง) วัตต์
\end{enumerate}
\item วงจร AC มี R = 30 Ω และ XC = 40 Ω (ไม่มี XL) อิมพีแดนซ์รวมเป็นเท่าไร \hfill \textbf{\small (Extra)}
\begin{enumerate}[label=)]
\item ก) 70 Ω
\item ข) 10 Ω
\item ค) 50 Ω
\item ง) 35 Ω
\end{enumerate}
\item cos φ ในสูตรกำลังไฟฟ้า AC เรียกอะไร \hfill \textbf{\small (Extra)}
\begin{enumerate}[label=)]
\item ก) มุมเฟส
\item ข) เพาเวอร์แฟกเตอร์
\item ค) รีแอกแตนซ์
\item ง) อิมพีแดนซ์
\end{enumerate}
\item วงจร AC มี P = 800 W, Vrms = 200 V และ Irms = 10 A cos φ มีค่าเท่าใด \hfill \textbf{\small (Extra)}
\begin{enumerate}[label=)]
\item ก) 0.2
\item ข) 0.4
\item ค) 0.8
\item ง) 1.0
\end{enumerate}
\item ตัวเหนี่ยวนำ L = 0.1 H ในวงจร AC ความถี่ 50 Hz มีค่า XL เท่าใด (π ≈ 3.14) \hfill \textbf{\small (Extra)}
\begin{enumerate}[label=)]
\item ก) 3.14 Ω
\item ข) 31.4 Ω
\item ค) 6.28 Ω
\item ง) 628 Ω
\end{enumerate}
\item ตัวเก็บประจุ C = 100 μF ในวงจร AC ความถี่ 50 Hz มีค่า XC เท่าใด (π ≈ 3.14) \hfill \textbf{\small (Extra)}
\begin{enumerate}[label=)]
\item ก) 3.14 Ω
\item ข) 31.8 Ω
\item ค) 318 Ω
\item ง) 6.28 Ω
\end{enumerate}
\item ข้อใดกล่าวถูกต้องเกี่ยวกับกระแสไฟฟ้าสลับ 50 Hz \hfill \textbf{\small (Extra)}
\begin{enumerate}[label=)]
\item ก) เปลี่ยนทิศทาง 25 ครั้งต่อวินาที
\item ข) เปลี่ยนทิศทาง 100 ครั้งต่อวินาที
\item ค) เปลี่ยนทิศทาง 50 ครั้งต่อวินาที
\item ง) เปลี่ยนทิศทาง 60 ครั้งต่อวินาที
\end{enumerate}
\item วงจร AC มี S = 1,000 VA และ cos φ = 0.6 กำลังไฟฟ้าจริง P เป็นเท่าไร \hfill \textbf{\small (Extra)}
\begin{enumerate}[label=)]
\item ก) 400 W
\item ข) 600 W
\item ค) 1,000 W
\item ง) 833 W
\end{enumerate}
\item วงจรอนุกรม RLC มี R = 100 Ω, XL = 200 Ω, XC = 50 Ω เมื่อความถี่เพิ่มขึ้น ข้อใดเกิดขึ้นกับ XC \hfill \textbf{\small (Extra)}
\begin{enumerate}[label=)]
\item ก) XC ลดลง
\item ข) XC เพิ่มขึ้น
\item ค) XC คงที่
\item ง) XC เป็นศูนย์
\end{enumerate}
\item วงจร AC มี R = 15 Ω และ XL = 20 Ω อิมพีแดนซ์เป็นเท่าไร \hfill \textbf{\small (Extra)}
\begin{enumerate}[label=)]
\item ก) 5 Ω
\item ข) 25 Ω
\item ค) 35 Ω
\item ง) 20 Ω
\end{enumerate}
\item หลอดไฟ 220 V/100 W ใช้กระแส RMS เท่าไรเมื่อทำงานที่แรงดันพิกัด \hfill \textbf{\small (Extra)}
\begin{enumerate}[label=)]
\item ก) 0.45 A
\item ข) 0.91 A
\item ค) 2.2 A
\item ง) 1 A
\end{enumerate}
\item วงจร AC มี R = 8 Ω, XL = 6 Ω, XC = 2 Ω ต่อกับแหล่งจ่าย 100 Vrms กำลังไฟฟ้าที่ใช้ไปมีค่าเท่าใด \hfill \textbf{\small (Extra)}
\begin{enumerate}[label=)]
\item ก) 800 W
\item ข) 500 W
\item ค) 640 W
\item ง) 1,000 W
\end{enumerate}
\item ข้อใดไม่ใช่ลักษณะของไฟฟ้ากระแสสลับ \hfill \textbf{\small (Extra)}
\begin{enumerate}[label=)]
\item ก) มีทั้งค่าบวกและค่าลบ
\item ข) มีความถี่คงที่
\item ค) มีค่าคงที่ตลอดเวลา
\item ง) สามารถแปลงแรงดันได้
\end{enumerate}
\item วงจรเรโซแนนซ์เกิดขึ้นเมื่อใด \hfill \textbf{\small (Extra)}
\begin{enumerate}[label=)]
\item ก) XL > XC
\item ข) XL < XC
\item ค) XL = XC
\item ง) R = 0
\end{enumerate}
\item วงจรเรโซแนนซ์มีเพาเวอร์แฟกเตอร์เท่าใด \hfill \textbf{\small (Extra)}
\begin{enumerate}[label=)]
\item ก) 0
\item ข) 0.5
\item ค) 1
\item ง) ขึ้นอยู่กับค่า R
\end{enumerate}
\item ตัวเก็บประจุ 50 μF ต่อกับแหล่งจ่าย AC ความถี่ 100 Hz มีค่า XC เท่าใด \hfill \textbf{\small (Extra)}
\begin{enumerate}[label=)]
\item ก) 3.18 Ω
\item ข) 31.8 Ω
\item ค) 318 Ω
\item ง) 6.28 Ω
\end{enumerate}
\item ความถี่เชิงมุม ω ของกระแส 50 Hz มีค่าเท่าไร (π ≈ 3.14) \hfill \textbf{\small (Extra)}
\begin{enumerate}[label=)]
\item ก) 100 rad/s
\item ข) 157 rad/s
\item ค) 314 rad/s
\item ง) 628 rad/s
\end{enumerate}
\item หม้อแปลงไฟฟ้าใช้หลักการใดทำงาน \hfill \textbf{\small (Extra)}
\begin{enumerate}[label=)]
\item ก) การเหนี่ยวนำแม่เหล็กไฟฟ้า
\item ข) การเหนี่ยวนำไฟฟ้าสถิต
\item ค) การแผ่รังสี
\item ง) การพาความร้อน
\end{enumerate}
\item แรงดัน V(t) = 300 sin(100πt) V มีค่า Vrms เท่าใด \hfill \textbf{\small (Extra)}
\begin{enumerate}[label=)]
\item ก) 150 V
\item ข) 212 V
\item ค) 300 V
\item ง) 424 V
\end{enumerate}
\item วงจรอนุกรม RLC มี L = 0.5 H, C = 100 μF, R = 10 Ω ที่ความถี่เรโซแนนซ์ ความถี่เท่าไร (π ≈ 3.14) \hfill \textbf{\small (Extra)}
\begin{enumerate}[label=)]
\item ก) ประมาณ 22.5 Hz
\item ข) ประมาณ 31.8 Hz
\item ค) ประมาณ 159 Hz
\item ง) ประมาณ 318 Hz
\end{enumerate}
\item กำลังไฟฟ้ารีแอกทีฟ Q มีหน่วยเป็นอะไร \hfill \textbf{\small (Extra)}
\begin{enumerate}[label=)]
\item ก) วัตต์ (W)
\item ข) โวลต์-แอมป์ (VA)
\item ค) VAR
\item ง) จูล (J)
\end{enumerate}
\item ไฟฟ้ากระแสตรง (DC) แตกต่างจากไฟฟ้ากระแสสลับ (AC) อย่างไร \hfill \textbf{\small (Extra)}
\begin{enumerate}[label=)]
\item ก) DC มีทิศทางเดียว AC เปลี่ยนทิศสม่ำเสมอ
\item ข) AC มีทิศทางเดียว DC เปลี่ยนทิศ
\item ค) DC เปลี่ยนทิศเร็วกว่า AC
\item ง) ไม่มีความแตกต่าง
\end{enumerate}
\item วงจร AC มี Vrms = 200 V, Z = 25 Ω และ cos φ = 0.6 กำลังไฟฟ้าเฉลี่ยมีค่าเท่าใด \hfill \textbf{\small (Extra)}
\begin{enumerate}[label=)]
\item ก) 800 W
\item ข) 960 W
\item ค) 1,600 W
\item ง) 480 W
\end{enumerate}
\item กระแสไฟฟ้าสลับ 220 Vrms มีค่าสูงสุดประมาณเท่าไร \hfill \textbf{\small (Extra)}
\begin{enumerate}[label=)]
\item ก) 110 V
\item ข) 156 V
\item ค) 220 V
\item ง) 311 V
\end{enumerate}
\item เตาไฟฟ้า AC มี cos φ = 0.85 กินกระแส 10 A ที่ 220 V กำลังไฟฟ้าเท่าไร \hfill \textbf{\small (Extra)}
\begin{enumerate}[label=)]
\item ก) 1,870 W
\item ข) 2,200 W
\item ค) 2,588 W
\item ง) 1,500 W
\end{enumerate}
\item วงจร AC มี R = 20 Ω, XL = 20 Ω และ XC = 20 Ω อิมพีแดนซ์และมุมเฟสเป็นเท่าไร \hfill \textbf{\small (Extra)}
\begin{enumerate}[label=)]
\item ก) Z = 20 Ω, φ = 0°
\item ข) Z = 20 Ω, φ = 45°
\item ค) Z = 40 Ω, φ = 0°
\item ง) Z = 20 Ω, φ = 90°
\end{enumerate}
\end{enumerate}
\end{document}
